\documentclass[11pt,a4paper]{article}
\usepackage[margin=2.5cm]{geometry}
\usepackage{amsmath,amssymb}
\usepackage{float}
\usepackage{booktabs}
\usepackage{multirow}
\usepackage{pdflscape}
\usepackage{array}
\usepackage{longtable}
\usepackage{hyperref}
\usepackage[T1]{fontenc}
\hypersetup{colorlinks=true,linkcolor=blue,citecolor=blue,urlcolor=blue}

\title{
  \textbf{Register Spill Analysis on RISC-V}\\[4pt]
  \large SPEC CPU2017 Benchmark Suite\\
  \normalsize gem5 TimingSimpleCPU $\cdot$ Custom SpillDetector $\cdot$ L1/L2 Spill Cache Counters
}
\author{
  L\"utfullah Burak Kaya\thanks{burak.kaya.50711@ozu.edu.tr, Ozyegin University}
  \and
  Ismail Akturk\thanks{ismail.akturk@ozyegin.edu.tr, Ozyegin University}
}
\date{\today}
\begin{document}
\maketitle
\tableofcontents
\clearpage


\section{Introduction}
\label{sec:intro}

Register spilling is the compiler-induced pattern of writing a live
register value to the stack (\emph{spill store}) and subsequently
reloading it (\emph{spill reload}) when the number of simultaneously
live values exceeds the ISA's physical register count.  On RISC-V
(RV64GC), the general-purpose and floating-point register files each
contain 32 entries.  Any computation that requires more than 32
live integer or floating-point values at once forces the compiler to
save the excess to memory.

This report presents the complete results of a custom gem5-based
spill measurement campaign covering all 15 C/C++ benchmarks from
SPEC CPU2017.  For each benchmark we report:
\begin{enumerate}
  \item ROI instruction, load, store, and spill counts for each
        available input tier (test, train, ref).
  \item Two spill-fraction metrics: \textbf{SP/inst\%}
        (spills\,/\,instructions) and \textbf{SP/load\%}
        (spills\,/\,loads).
  \item For the test input and for ref runs performed with
        \texttt{--l2cache}: L1-D and L2 cache hit/miss rates
        for spill loads and spill stores separately.
\end{enumerate}

Section~\ref{sec:method} describes the measurement infrastructure in
detail, including the rationale for each metric.
Sections~\ref{sec:deepsjeng}--\ref{sec:perlbench} present
per-benchmark results.
Section~\ref{sec:summary} contains a cross-benchmark comparison.


\section{Methodology}
\label{sec:method}

\subsection{gem5 Simulation Environment}

All simulations use gem5 Syscall-Emulation (SE) mode with:
\begin{itemize}
  \item \textbf{CPU model}: TimingSimpleCPU (in-order, timing-accurate)
  \item \textbf{ISA}: RISC-V RV64GC
  \item \textbf{L1-D cache}: 32\,kB, 8-way, 64\,B lines
  \item \textbf{L2 cache}: 256\,kB, 8-way unified, write-back
  \item \textbf{DRAM}: 4\,GB SimpleMemory ($\approx$63\,ns latency)
  \item \textbf{gem5 flags}: \texttt{--cpu-type=TimingSimpleCPU --caches --l2cache --mem-size=4GB}
\end{itemize}

\subsection{SpillDetector Module}

A custom \texttt{SpillDetector} is integrated into
\texttt{gem5/src/cpu/simple/timing.cc}.  It operates on two hooks:
\begin{enumerate}
  \item \textbf{onStoreInstruction}: every store instruction inside the
    ROI is checked.  If the target address is a \emph{stack address}
    (within the current SP $\pm$ \texttt{MAX\_SPILL\_SIZE} bytes), the
    address and issue tick are recorded in a hash map.
  \item \textbf{onLoadInstruction}: every load inside the ROI is checked.
    Stale entries older than \texttt{MAX\_SPILL\_WINDOW}
    (50 million ticks $\approx$ 50\,\textmu s) are first purged
    to prevent false positives from long-lived stack data.
    Then, if the load address matches a recorded store entry, it is
    counted as a \emph{spill reload} (\texttt{roi\_spills}).
\end{enumerate}

The detector outputs a per-run file
\texttt{riscv\_spill\_stats.txt} containing
\texttt{roi\_instructions}, \texttt{roi\_loads}, \texttt{roi\_stores},
and \texttt{roi\_spills}.

\subsection{Spill-Tagged Cache Counters}

Two new \texttt{Request} flags---\texttt{SPILL\_LOAD} and
\texttt{SPILL\_STORE}---are set on memory requests that the
SpillDetector identifies as spill operations.  The L1-D and L2
caches separately accumulate hit and miss counters for
spill-tagged requests, producing four new per-cache statistics:
\texttt{spillLoadHits}, \texttt{spillLoadMisses},
\texttt{spillStoreHits}, \texttt{spillStoreMisses}.

\subsection{ROI Markers}

Each benchmark binary was compiled with gem5 ROI markers
(\texttt{m5\_work\_begin()} / \texttt{m5\_work\_end()}) around
the dominant computational kernel, as listed in Table~\ref{tab:roi_markers}.

\begin{table}[H]
\centering
\caption{ROI marker placement per benchmark}
\label{tab:roi_markers}
\small
\begin{tabular}{lll}
\toprule
\textbf{Benchmark} & \textbf{Source file} & \textbf{ROI wraps} \\
\midrule
531.deepsjeng\_r & \texttt{src/sjeng.cpp}            & \texttt{run\_epd\_testsuite()} \\
541.leela\_r     & \texttt{src/Leela.cpp}             & main try/catch block \\
508.namd\_r      & \texttt{src/spec\_namd.C}          & iterations for-loop \\
519.lbm\_r       & \texttt{src/main.c}                & time-step for-loop \\
544.nab\_r       & \texttt{src/nabmd.c}               & two \texttt{md()} calls \\
557.xz\_r        & \texttt{src/spec.c}                & compression for-loop \\
538.imagick\_r   & \texttt{src/utilities/convert.c}   & \texttt{ConvertMain()} \\
526.blender\_r   & \texttt{...cycles\_standalone.cpp} & session init + wait \\
511.povray\_r    & \texttt{src/povray.cpp}            & \texttt{StartRender()} loop \\
520.omnetpp\_r   & \texttt{src/simulator/main.cc}     & \texttt{setupUserInterface()} \\
523.xalancbmk\_r & \texttt{src/XalanExe.cpp}          & \texttt{xsltMain()} \\
525.x264\_r      & \texttt{src/ldecod\_src/...}       & do-while decode loop \\
502.gcc\_r       & \texttt{src/main.c}                & \texttt{toplev\_main()} \\
505.mcf\_r       & \texttt{src/mcf.c}                 & \texttt{global\_opt()} \\
500.perlbench\_r & \texttt{src/perlmain.c}            & \texttt{perl\_run()} \\
\bottomrule
\end{tabular}
\end{table}

\subsection{Metric Definitions and Rationale}
\label{subsec:metrics}

Let $N_L = \textit{spillLoadHits}_\text{L1} + \textit{spillLoadMisses}_\text{L1}$ and
$N_S = \textit{spillStoreHits}_\text{L1} + \textit{spillStoreMisses}_\text{L1}$.
All cache-level percentages share the same denominator so that
$\text{L1H\%}+\text{L2H\%}+\text{DRAML\%} = 100\%$
(and identically for stores).

\begin{align}
\text{SP/inst\%}  &= \frac{\text{roi\_spills}}{\text{roi\_instructions}} \times 100
\label{eq:spinst}\\
\text{SP/load\%}  &= \frac{\text{roi\_spills}}{\text{roi\_loads}} \times 100
\label{eq:spload}\\
\text{L1H\%}  &= \frac{\textit{spillLoadHits}_\text{L1}}{N_L} \times 100
\label{eq:l1h}\\
\text{L2H\%}  &= \frac{\textit{spillLoadHits}_\text{L2}}{N_L} \times 100
\label{eq:l2h}\\
\text{DRAML\%}&= \frac{\textit{spillLoadMisses}_\text{L2}}{N_L} \times 100
\label{eq:draml}
\end{align}
L1HS\%, L2HS\%, and DRAMS\% use the same formulae with spill-store counters.

\paragraph{Why SP/inst\%?}
Equation~\eqref{eq:spinst} measures the \emph{density} of spill
operations within the instruction stream.  It answers ``what fraction
of the program's issue bandwidth is consumed by register-spill memory
traffic?''  A value of 4\% means 1 in every 25 instructions is a
spill reload.  This is the most direct proxy for the \emph{code-size
and issue-slot overhead} of spilling.

\paragraph{Why SP/load\%?}
Equation~\eqref{eq:spload} measures the \emph{fraction of load traffic}
attributable to spilling.  It is arguably the more important metric
for memory-system analysis because the load unit and cache are shared
resources: displacing useful (non-spill) loads by spill reloads
wastes bandwidth.  A value of 30\% means that eliminating spilling
would reduce load traffic by 30\%, with direct impact on CPI and
energy.

\paragraph{Why L1H\% and L2H\%?}
Equations~\eqref{eq:l1h}--\eqref{eq:l2h} characterise the
\emph{cache footprint} of spilled data.  Because a spill store and
its subsequent spill reload are temporally adjacent (the value is
reloaded as soon as a register is needed again), the compiler's stack
discipline naturally produces high spatial and temporal locality.
We expect---and confirm---that spill traffic is overwhelmingly served
by the L1-D cache.  If L1H\% were low, it would suggest that the
ROI involves many long-lived variables that are pushed beyond L1's
capacity before being reloaded, which would indicate a qualitatively
different spilling regime.

\paragraph{Why DRAML\%?}
Equation~\eqref{eq:draml} quantifies whether spilling imposes
\emph{off-chip memory traffic}.  Near-zero DRAML\% (as observed
across the suite) means that spilling costs latency and issue slots
but does \emph{not} burden the memory bus.  This is an important
distinction: a hardware register-file extension that eliminated
spilling would primarily recover CPU cycles rather than DRAM bandwidth.

\paragraph{Why L2H\% alongside DRAML\%?}
L2H\% measures what fraction of \emph{all} spill loads are served by
L2 (not a conditional rate on L1 misses).  Together L1H\%+L2H\%+DRAML\%$=$100\%,
so L2H\% quantifies the share of total spill-load traffic that hits L2.
Near-zero DRAML\% with L2H\%$\approx$L1M\% confirms that L2 rescues
every L1 spill miss with no escapes to DRAM.

\clearpage


%%──────────────────────────────────────────────────────────────────────────────
\section{\texttt{531.deepsjeng\_r}}
\label{sec:deepsjeng}

\paragraph{Benchmark summary.}
\texttt{531.deepsjeng\_r} is the SPEC CPU2017 rate variant of Deep Sjeng.
Recursive alpha-beta game-tree search with large call stacks and transposition table bookkeeping.
The ROI is bounded by \texttt{m5\_work\_begin()} / \texttt{m5\_work\_end()}
around \texttt{run\_epd\_testsuite()}.

\subsection{Input Datasets}

\begin{table}[H]
\centering
\caption{Input dataset sizes for \texttt{531.deepsjeng\_r}}
\label{tab:deepsjeng_inputs}
\small
\begin{tabular}{lp{8cm}l}
\toprule
\textbf{Tier} & \textbf{Input} & \textbf{Relative scale} \\
\midrule
\textbf{test} & test.txt --- 2 chess positions, search depth 15 & 2 positions \\
train & train.txt --- 20 positions, search depth 15 & 10$\times$ more positions \\
ref & ref.txt --- 24 positions, search depth 15 & 12$\times$ more positions \\
\bottomrule
\end{tabular}
\end{table}

\subsection{Test Input Results}

\begin{table}[H]
\centering
\caption{ROI spill statistics --- \texttt{531.deepsjeng\_r} (test input)}
\label{tab:deepsjeng_test_roi}
\begin{tabular}{lr}
\toprule
\textbf{Metric} & \textbf{Value} \\
\midrule
ROI instructions & 46,202,686,170 \\
ROI loads        & 7,989,081,366 \\
ROI stores       & 3,395,185,975 \\
\textbf{ROI spills}   & \textbf{1,975,804,388} \\
\textbf{SP/inst\%}   & \textbf{4.28\%} \\
\textbf{SP/load\%}   & \textbf{24.73\%} \\
Simulated time   & 74.082\,s \\
CPI              & 3.21 \\
\bottomrule
\end{tabular}
\end{table}

\begin{table}[H]
\centering
\caption{Spill cache hit/miss breakdown --- \texttt{531.deepsjeng\_r} (test input)}
\label{tab:deepsjeng_test_cache}
\begin{tabular}{lrrrrr}
\toprule
\textbf{} & \textbf{Total (M)} & \textbf{L1H\%} & \textbf{L1M\%} & \textbf{L2H\%} & \textbf{DRAM\%} \\
\midrule
Spill loads  & 1,975.8 & 99.97 & 0.03 & 0.03 & 0.00 \\
Spill stores & 2,462.9 & 99.68 & 0.32 & 0.32 & 0.00 \\
\multicolumn{2}{l}{\small\textit{L1H\%+L2H\%+DRAM\%=100\%}} \\
\bottomrule
\end{tabular}
\end{table}

\subsection{Train Input Results}

\textit{Note: this run used an older simulation without \texttt{-{}-l2cache}; cache counters are not available.}


\begin{table}[H]
\centering
\caption{ROI spill statistics --- \texttt{531.deepsjeng\_r} (train input)}
\label{tab:deepsjeng_train_roi}
\begin{tabular}{lr}
\toprule
\textbf{Metric} & \textbf{Value} \\
\midrule
ROI instructions & 399,987,651,204 \\
ROI loads        & 72,368,911,103 \\
ROI stores       & 32,213,674,285 \\
\textbf{ROI spills}   & \textbf{18,131,499,845} \\
\textbf{SP/inst\%}   & \textbf{4.53\%} \\
\textbf{SP/load\%}   & \textbf{25.05\%} \\
\bottomrule
\end{tabular}
\end{table}

\subsection{Ref Input Results}
\textit{Ref simulation in progress.}



%%──────────────────────────────────────────────────────────────────────────────
\section{\texttt{541.leela\_r}}
\label{sec:leela}

\paragraph{Benchmark summary.}
\texttt{541.leela\_r} is the SPEC CPU2017 rate variant of Leela.
Monte-Carlo tree search with per-node statistics; high function-call depth drives register pressure.
The ROI is bounded by \texttt{m5\_work\_begin()} / \texttt{m5\_work\_end()}
around \texttt{main try/catch block}.

\subsection{Input Datasets}

\begin{table}[H]
\centering
\caption{Input dataset sizes for \texttt{541.leela\_r}}
\label{tab:leela_inputs}
\small
\begin{tabular}{lp{8cm}l}
\toprule
\textbf{Tier} & \textbf{Input} & \textbf{Relative scale} \\
\midrule
\textbf{test} & test.sgf --- 1.6\,kB Go game record & 1.6\,kB \\
train & train.sgf --- 642\,B (shorter game) & shorter than test \\
ref & ref.sgf --- 4.7\,kB (longer game) & 3$\times$ larger than test \\
\bottomrule
\end{tabular}
\end{table}

\subsection{Test Input Results}

\begin{table}[H]
\centering
\caption{ROI spill statistics --- \texttt{541.leela\_r} (test input)}
\label{tab:leela_test_roi}
\begin{tabular}{lr}
\toprule
\textbf{Metric} & \textbf{Value} \\
\midrule
ROI instructions & 26,063,891,645 \\
ROI loads        & 5,110,627,713 \\
ROI stores       & 1,756,742,699 \\
\textbf{ROI spills}   & \textbf{1,179,370,892} \\
\textbf{SP/inst\%}   & \textbf{4.52\%} \\
\textbf{SP/load\%}   & \textbf{23.08\%} \\
Simulated time   & 43.332\,s \\
CPI              & 3.33 \\
\bottomrule
\end{tabular}
\end{table}

\begin{table}[H]
\centering
\caption{Spill cache hit/miss breakdown --- \texttt{541.leela\_r} (test input)}
\label{tab:leela_test_cache}
\begin{tabular}{lrrrrr}
\toprule
\textbf{} & \textbf{Total (M)} & \textbf{L1H\%} & \textbf{L1M\%} & \textbf{L2H\%} & \textbf{DRAM\%} \\
\midrule
Spill loads  & 1,179.4 & 100.00 & 0.00 & 0.00 & 0.00 \\
Spill stores & 1,430.8 & 99.81 & 0.19 & 0.19 & 0.00 \\
\multicolumn{2}{l}{\small\textit{L1H\%+L2H\%+DRAM\%=100\%}} \\
\bottomrule
\end{tabular}
\end{table}

\subsection{Train Input Results}

\textit{Note: this run used an older simulation without \texttt{-{}-l2cache}; cache counters are not available.}


\begin{table}[H]
\centering
\caption{ROI spill statistics --- \texttt{541.leela\_r} (train input)}
\label{tab:leela_train_roi}
\begin{tabular}{lr}
\toprule
\textbf{Metric} & \textbf{Value} \\
\midrule
ROI instructions & 436,316,211,710 \\
ROI loads        & 85,302,835,628 \\
ROI stores       & 27,570,532,075 \\
\textbf{ROI spills}   & \textbf{19,460,747,256} \\
\textbf{SP/inst\%}   & \textbf{4.46\%} \\
\textbf{SP/load\%}   & \textbf{22.81\%} \\
\bottomrule
\end{tabular}
\end{table}

\subsection{Ref Input Results}
\textit{Ref simulation in progress.}



%%──────────────────────────────────────────────────────────────────────────────
\section{\texttt{508.namd\_r}}
\label{sec:namd}

\paragraph{Benchmark summary.}
\texttt{508.namd\_r} is the SPEC CPU2017 rate variant of NAMD.
Vectorised molecular-dynamics force computation; SIMD-friendly loops with moderate spill rate.
The ROI is bounded by \texttt{m5\_work\_begin()} / \texttt{m5\_work\_end()}
around \texttt{iterations for-loop in spec\_namd.C}.

\subsection{Input Datasets}

\begin{table}[H]
\centering
\caption{Input dataset sizes for \texttt{508.namd\_r}}
\label{tab:namd_inputs}
\small
\begin{tabular}{lp{8cm}l}
\toprule
\textbf{Tier} & \textbf{Input} & \textbf{Relative scale} \\
\midrule
\textbf{test} & apoa1.input, \textbf{1 iteration} & 1 iter \\
train & apoa1.input, \textbf{7 iterations} & 7$\times$ iterations \\
ref & apoa1.input, \textbf{65 iterations} & 65$\times$ iterations \\
\bottomrule
\end{tabular}
\end{table}

\subsection{Test Input Results}

\begin{table}[H]
\centering
\caption{ROI spill statistics --- \texttt{508.namd\_r} (test input)}
\label{tab:namd_test_roi}
\begin{tabular}{lr}
\toprule
\textbf{Metric} & \textbf{Value} \\
\midrule
ROI instructions & 31,914,751,100 \\
ROI loads        & 8,380,476,926 \\
ROI stores       & 1,945,521,360 \\
\textbf{ROI spills}   & \textbf{335,693,695} \\
\textbf{SP/inst\%}   & \textbf{1.05\%} \\
\textbf{SP/load\%}   & \textbf{4.01\%} \\
Simulated time   & 60.297\,s \\
CPI              & 3.78 \\
\bottomrule
\end{tabular}
\end{table}

\begin{table}[H]
\centering
\caption{Spill cache hit/miss breakdown --- \texttt{508.namd\_r} (test input)}
\label{tab:namd_test_cache}
\begin{tabular}{lrrrrr}
\toprule
\textbf{} & \textbf{Total (M)} & \textbf{L1H\%} & \textbf{L1M\%} & \textbf{L2H\%} & \textbf{DRAM\%} \\
\midrule
Spill loads  & 335.7 & 99.93 & 0.07 & 0.07 & 0.00 \\
Spill stores & 345.7 & 99.82 & 0.18 & 0.18 & 0.00 \\
\multicolumn{2}{l}{\small\textit{L1H\%+L2H\%+DRAM\%=100\%}} \\
\bottomrule
\end{tabular}
\end{table}

\subsection{Train Input Results}

\textit{Note: this run used an older simulation without \texttt{-{}-l2cache}; cache counters are not available.}


\begin{table}[H]
\centering
\caption{ROI spill statistics --- \texttt{508.namd\_r} (train input)}
\label{tab:namd_train_roi}
\begin{tabular}{lr}
\toprule
\textbf{Metric} & \textbf{Value} \\
\midrule
ROI instructions & 218,427,292,249 \\
ROI loads        & 57,662,937,316 \\
ROI stores       & 13,245,362,832 \\
\textbf{ROI spills}   & \textbf{2,285,717,035} \\
\textbf{SP/inst\%}   & \textbf{1.05\%} \\
\textbf{SP/load\%}   & \textbf{3.96\%} \\
\bottomrule
\end{tabular}
\end{table}

\subsection{Ref Input Results}
\textit{Ref simulation in progress.}



%%──────────────────────────────────────────────────────────────────────────────
\section{\texttt{519.lbm\_r}}
\label{sec:lbm}

\paragraph{Benchmark summary.}
\texttt{519.lbm\_r} is the SPEC CPU2017 rate variant of LBM.
Single innermost kernel sweeping a 3-D lattice; 19 FP distribution values per cell exceed the 32 available FP registers.
The ROI is bounded by \texttt{m5\_work\_begin()} / \texttt{m5\_work\_end()}
around \texttt{main time-step for-loop}.

\subsection{Input Datasets}

\begin{table}[H]
\centering
\caption{Input dataset sizes for \texttt{519.lbm\_r}}
\label{tab:lbm_inputs}
\small
\begin{tabular}{lp{8cm}l}
\toprule
\textbf{Tier} & \textbf{Input} & \textbf{Relative scale} \\
\midrule
\textbf{test} & \textbf{20} time-steps, channel-flow domain (cf\_a.of) & 20 steps \\
train & \textbf{300} time-steps, channel-flow domain (cf\_b.of) & 15$\times$ more steps \\
ref & \textbf{3000} time-steps, lid-driven cavity (ldc.of) & 150$\times$ more steps \\
\bottomrule
\end{tabular}
\end{table}

\subsection{Test Input Results}

\begin{table}[H]
\centering
\caption{ROI spill statistics --- \texttt{519.lbm\_r} (test input)}
\label{tab:lbm_test_roi}
\begin{tabular}{lr}
\toprule
\textbf{Metric} & \textbf{Value} \\
\midrule
ROI instructions & 6,506,800,041 \\
ROI loads        & 1,276,836,306 \\
ROI stores       & 775,731,184 \\
\textbf{ROI spills}   & \textbf{274,130,440} \\
\textbf{SP/inst\%}   & \textbf{4.21\%} \\
\textbf{SP/load\%}   & \textbf{21.47\%} \\
Simulated time   & 35.835\,s \\
CPI              & 11.01 \\
\bottomrule
\end{tabular}
\end{table}

\begin{table}[H]
\centering
\caption{Spill cache hit/miss breakdown --- \texttt{519.lbm\_r} (test input)}
\label{tab:lbm_test_cache}
\begin{tabular}{lrrrrr}
\toprule
\textbf{} & \textbf{Total (M)} & \textbf{L1H\%} & \textbf{L1M\%} & \textbf{L2H\%} & \textbf{DRAM\%} \\
\midrule
Spill loads  & 274.1 & 100.00 & 0.00 & 0.00 & 0.00 \\
Spill stores & 274.1 & 99.92 & 0.08 & 0.08 & 0.00 \\
\multicolumn{2}{l}{\small\textit{L1H\%+L2H\%+DRAM\%=100\%}} \\
\bottomrule
\end{tabular}
\end{table}

\subsection{Train Input Results}

\textit{Note: this run used an older simulation without \texttt{-{}-l2cache}; cache counters are not available.}


\begin{table}[H]
\centering
\caption{ROI spill statistics --- \texttt{519.lbm\_r} (train input)}
\label{tab:lbm_train_roi}
\begin{tabular}{lr}
\toprule
\textbf{Metric} & \textbf{Value} \\
\midrule
ROI instructions & 100,160,472,569 \\
ROI loads        & 19,574,976,720 \\
ROI stores       & 11,753,652,015 \\
\textbf{ROI spills}   & \textbf{4,229,639,284} \\
\textbf{SP/inst\%}   & \textbf{4.22\%} \\
\textbf{SP/load\%}   & \textbf{21.61\%} \\
\bottomrule
\end{tabular}
\end{table}

\subsection{Ref Input Results}
\textit{Ref simulation in progress.}



%%──────────────────────────────────────────────────────────────────────────────
\section{\texttt{544.nab\_r}}
\label{sec:nab}

\paragraph{Benchmark summary.}
\texttt{544.nab\_r} is the SPEC CPU2017 rate variant of NAB.
Implicit-solvent MD with Born and vacuum phases; two separate \texttt{md()} calls share the ROI.
The ROI is bounded by \texttt{m5\_work\_begin()} / \texttt{m5\_work\_end()}
around \texttt{two md() calls (Born + vacuum)}.

\subsection{Input Datasets}

\begin{table}[H]
\centering
\caption{Input dataset sizes for \texttt{544.nab\_r}}
\label{tab:nab_inputs}
\small
\begin{tabular}{lp{8cm}l}
\toprule
\textbf{Tier} & \textbf{Input} & \textbf{Relative scale} \\
\midrule
\textbf{test} & hkrdenq molecule, 1000 MD steps & small molecule \\
train & aminos molecule, 1000 MD steps & larger molecule \\
ref & 1am0 molecule, 122 MD steps & largest molecule \\
\bottomrule
\end{tabular}
\end{table}

\subsection{Test Input Results}

\begin{table}[H]
\centering
\caption{ROI spill statistics --- \texttt{544.nab\_r} (test input)}
\label{tab:nab_test_roi}
\begin{tabular}{lr}
\toprule
\textbf{Metric} & \textbf{Value} \\
\midrule
ROI instructions & 6,691,483,052 \\
ROI loads        & 1,523,059,851 \\
ROI stores       & 321,568,363 \\
\textbf{ROI spills}   & \textbf{153,736,211} \\
\textbf{SP/inst\%}   & \textbf{2.30\%} \\
\textbf{SP/load\%}   & \textbf{10.09\%} \\
Simulated time   & 10.840\,s \\
CPI              & 3.24 \\
\bottomrule
\end{tabular}
\end{table}

\begin{table}[H]
\centering
\caption{Spill cache hit/miss breakdown --- \texttt{544.nab\_r} (test input)}
\label{tab:nab_test_cache}
\begin{tabular}{lrrrrr}
\toprule
\textbf{} & \textbf{Total (M)} & \textbf{L1H\%} & \textbf{L1M\%} & \textbf{L2H\%} & \textbf{DRAM\%} \\
\midrule
Spill loads  & 153.7 & 100.00 & 0.00 & 0.00 & 0.00 \\
Spill stores & 159.2 & 99.99 & 0.01 & 0.01 & 0.00 \\
\multicolumn{2}{l}{\small\textit{L1H\%+L2H\%+DRAM\%=100\%}} \\
\bottomrule
\end{tabular}
\end{table}

\subsection{Train Input Results}

\textit{Train simulation in progress.}


\subsection{Ref Input Results}
\textit{Ref simulation in progress.}



%%──────────────────────────────────────────────────────────────────────────────
\section{\texttt{557.xz\_r}}
\label{sec:xz}

\paragraph{Benchmark summary.}
\texttt{557.xz\_r} is the SPEC CPU2017 rate variant of XZ.
LZMA2 encoder/decoder with large sliding-window state; compression inner loops sustain register pressure.
The ROI is bounded by \texttt{m5\_work\_begin()} / \texttt{m5\_work\_end()}
around \texttt{compression for-loop in spec.c}.

\subsection{Input Datasets}

\begin{table}[H]
\centering
\caption{Input dataset sizes for \texttt{557.xz\_r}}
\label{tab:xz_inputs}
\small
\begin{tabular}{lp{8cm}l}
\toprule
\textbf{Tier} & \textbf{Input} & \textbf{Relative scale} \\
\midrule
\textbf{test} & cpu2006docs.tar.xz --- 1.3\,MB, compression level 0 & 1.3\,MB file \\
train & IMG\_2560.cr2.xz --- 15\,MB, compression level 4 & 15\,MB (12$\times$), harder \\
ref & cld.tar.xz --- 79\,MB, compression level 6 & 79\,MB (61$\times$), hardest \\
\bottomrule
\end{tabular}
\end{table}

\subsection{Test Input Results}

\begin{table}[H]
\centering
\caption{ROI spill statistics --- \texttt{557.xz\_r} (test input)}
\label{tab:xz_test_roi}
\begin{tabular}{lr}
\toprule
\textbf{Metric} & \textbf{Value} \\
\midrule
ROI instructions & 2,014,351,522 \\
ROI loads        & 367,733,280 \\
ROI stores       & 183,548,298 \\
\textbf{ROI spills}   & \textbf{39,181,468} \\
\textbf{SP/inst\%}   & \textbf{1.95\%} \\
\textbf{SP/load\%}   & \textbf{10.65\%} \\
Simulated time   & 4.802\,s \\
CPI              & 4.77 \\
\bottomrule
\end{tabular}
\end{table}

\begin{table}[H]
\centering
\caption{Spill cache hit/miss breakdown --- \texttt{557.xz\_r} (test input)}
\label{tab:xz_test_cache}
\begin{tabular}{lrrrrr}
\toprule
\textbf{} & \textbf{Total (M)} & \textbf{L1H\%} & \textbf{L1M\%} & \textbf{L2H\%} & \textbf{DRAM\%} \\
\midrule
Spill loads  & 39.2 & 100.00 & 0.00 & 0.00 & 0.00 \\
Spill stores & 42.6 & 99.90 & 0.10 & 0.10 & 0.00 \\
\multicolumn{2}{l}{\small\textit{L1H\%+L2H\%+DRAM\%=100\%}} \\
\bottomrule
\end{tabular}
\end{table}

\subsection{Train Input Results}

\textit{Note: this run used an older simulation without \texttt{-{}-l2cache}; cache counters are not available.}


\begin{table}[H]
\centering
\caption{ROI spill statistics --- \texttt{557.xz\_r} (train input)}
\label{tab:xz_train_roi}
\begin{tabular}{lr}
\toprule
\textbf{Metric} & \textbf{Value} \\
\midrule
ROI instructions & 45,071,901,346 \\
ROI loads        & 8,781,014,341 \\
ROI stores       & 5,022,533,854 \\
\textbf{ROI spills}   & \textbf{1,422,403,651} \\
\textbf{SP/inst\%}   & \textbf{3.16\%} \\
\textbf{SP/load\%}   & \textbf{16.20\%} \\
\bottomrule
\end{tabular}
\end{table}

\subsection{Ref Input Results}

\textit{Note: this run used the older simulation without \texttt{-{}-l2cache}; L1/L2 cache counters are not available.}


\begin{table}[H]
\centering
\caption{ROI spill statistics --- \texttt{557.xz\_r} (ref input)}
\label{tab:xz_ref_roi}
\begin{tabular}{lr}
\toprule
\textbf{Metric} & \textbf{Value} \\
\midrule
ROI instructions & 401,325,609,454 \\
ROI loads        & 65,376,336,569 \\
ROI stores       & 24,760,673,373 \\
\textbf{ROI spills}   & \textbf{6,636,190,731} \\
\textbf{SP/inst\%}   & \textbf{1.65\%} \\
\textbf{SP/load\%}   & \textbf{10.15\%} \\
\bottomrule
\end{tabular}
\end{table}


%%──────────────────────────────────────────────────────────────────────────────
\section{\texttt{538.imagick\_r}}
\label{sec:imagick}

\paragraph{Benchmark summary.}
\texttt{538.imagick\_r} is the SPEC CPU2017 rate variant of ImageMagick.
Image-processing pipeline with chained convolution filters; the tiny test image means very short ROI.
The ROI is bounded by \texttt{m5\_work\_begin()} / \texttt{m5\_work\_end()}
around \texttt{ConvertMain()}.

\subsection{Input Datasets}

\begin{table}[H]
\centering
\caption{Input dataset sizes for \texttt{538.imagick\_r}}
\label{tab:imagick_inputs}
\small
\begin{tabular}{lp{8cm}l}
\toprule
\textbf{Tier} & \textbf{Input} & \textbf{Relative scale} \\
\midrule
\textbf{test} & test\_input.tga (318\,B), 4-filter chain & tiny synthetic image \\
train & train\_input.tga (318\,B), 8-filter chain with resize & same image, more filters \\
ref & refrate\_input.tga (8.2\,MB), 6-filter chain & 8.2\,MB real image \\
\bottomrule
\end{tabular}
\end{table}

\subsection{Test Input Results}

\begin{table}[H]
\centering
\caption{ROI spill statistics --- \texttt{538.imagick\_r} (test input)}
\label{tab:imagick_test_roi}
\begin{tabular}{lr}
\toprule
\textbf{Metric} & \textbf{Value} \\
\midrule
ROI instructions & 117,707,681 \\
ROI loads        & 20,608,164 \\
ROI stores       & 8,985,699 \\
\textbf{ROI spills}   & \textbf{2,426,438} \\
\textbf{SP/inst\%}   & \textbf{2.06\%} \\
\textbf{SP/load\%}   & \textbf{11.77\%} \\
Simulated time   & 0.205\,s \\
CPI              & 3.49 \\
\bottomrule
\end{tabular}
\end{table}

\begin{table}[H]
\centering
\caption{Spill cache hit/miss breakdown --- \texttt{538.imagick\_r} (test input)}
\label{tab:imagick_test_cache}
\begin{tabular}{lrrrrr}
\toprule
\textbf{} & \textbf{Total (M)} & \textbf{L1H\%} & \textbf{L1M\%} & \textbf{L2H\%} & \textbf{DRAM\%} \\
\midrule
Spill loads  & 2.4 & 100.00 & 0.00 & 0.00 & 0.00 \\
Spill stores & 2.5 & 99.74 & 0.26 & 0.22 & 0.04 \\
\multicolumn{2}{l}{\small\textit{L1H\%+L2H\%+DRAM\%=100\%}} \\
\bottomrule
\end{tabular}
\end{table}

\subsection{Train Input Results}

\textit{Note: this run used an older simulation without \texttt{-{}-l2cache}; cache counters are not available.}


\begin{table}[H]
\centering
\caption{ROI spill statistics --- \texttt{538.imagick\_r} (train input)}
\label{tab:imagick_train_roi}
\begin{tabular}{lr}
\toprule
\textbf{Metric} & \textbf{Value} \\
\midrule
ROI instructions & 284,114,913,819 \\
ROI loads        & 47,764,394,626 \\
ROI stores       & 652,850,800 \\
\textbf{ROI spills}   & \textbf{292,429,548} \\
\textbf{SP/inst\%}   & \textbf{0.10\%} \\
\textbf{SP/load\%}   & \textbf{0.61\%} \\
\bottomrule
\end{tabular}
\end{table}

\subsection{Ref Input Results}

\begin{table}[H]
\centering
\caption{ROI spill statistics --- \texttt{538.imagick\_r} (ref input)}
\label{tab:imagick_ref_roi}
\begin{tabular}{lr}
\toprule
\textbf{Metric} & \textbf{Value} \\
\midrule
ROI instructions & 117,706,531 \\
ROI loads        & 20,607,957 \\
ROI stores       & 8,985,593 \\
\textbf{ROI spills}   & \textbf{2,426,383} \\
\textbf{SP/inst\%}   & \textbf{2.06\%} \\
\textbf{SP/load\%}   & \textbf{11.77\%} \\
Simulated time   & 0.205\,s \\
CPI              & 3.49 \\
\bottomrule
\end{tabular}
\end{table}

\begin{table}[H]
\centering
\caption{Spill cache hit/miss breakdown --- \texttt{538.imagick\_r} (ref input)}
\label{tab:imagick_ref_cache}
\begin{tabular}{lrrrrr}
\toprule
\textbf{} & \textbf{Total (M)} & \textbf{L1H\%} & \textbf{L1M\%} & \textbf{L2H\%} & \textbf{DRAM\%} \\
\midrule
Spill loads  & 2.4 & 100.00 & 0.00 & 0.00 & 0.00 \\
Spill stores & 2.5 & 99.74 & 0.26 & 0.22 & 0.04 \\
\multicolumn{2}{l}{\small\textit{L1H\%+L2H\%+DRAM\%=100\%}} \\
\bottomrule
\end{tabular}
\end{table}


%%──────────────────────────────────────────────────────────────────────────────
\section{\texttt{526.blender\_r}}
\label{sec:blender}

\paragraph{Benchmark summary.}
\texttt{526.blender\_r} is the SPEC CPU2017 rate variant of Blender.
Cycles path-tracer with recursive BVH traversal, shading, and per-sample RNG state -- second-highest spill rate.
The ROI is bounded by \texttt{m5\_work\_begin()} / \texttt{m5\_work\_end()}
around \texttt{session\_init() + wait + exit}.

\subsection{Input Datasets}

\begin{table}[H]
\centering
\caption{Input dataset sizes for \texttt{526.blender\_r}}
\label{tab:blender_inputs}
\small
\begin{tabular}{lp{8cm}l}
\toprule
\textbf{Tier} & \textbf{Input} & \textbf{Relative scale} \\
\midrule
\textbf{test} & cube.blend --- simple cube scene, frame 1 & trivial scene \\
train & sh5\_reduced.blend --- scene with objects, frame 234 & moderate complexity \\
ref & sh3\_no\_char.blend --- full production scene, frame 849 & full complexity \\
\bottomrule
\end{tabular}
\end{table}

\subsection{Test Input Results}

\begin{table}[H]
\centering
\caption{ROI spill statistics --- \texttt{526.blender\_r} (test input)}
\label{tab:blender_test_roi}
\begin{tabular}{lr}
\toprule
\textbf{Metric} & \textbf{Value} \\
\midrule
ROI instructions & 994,414,403 \\
ROI loads        & 233,612,260 \\
ROI stores       & 116,509,608 \\
\textbf{ROI spills}   & \textbf{64,831,836} \\
\textbf{SP/inst\%}   & \textbf{6.52\%} \\
\textbf{SP/load\%}   & \textbf{27.75\%} \\
Simulated time   & 2.076\,s \\
CPI              & 4.18 \\
\bottomrule
\end{tabular}
\end{table}

\begin{table}[H]
\centering
\caption{Spill cache hit/miss breakdown --- \texttt{526.blender\_r} (test input)}
\label{tab:blender_test_cache}
\begin{tabular}{lrrrrr}
\toprule
\textbf{} & \textbf{Total (M)} & \textbf{L1H\%} & \textbf{L1M\%} & \textbf{L2H\%} & \textbf{DRAM\%} \\
\midrule
Spill loads  & 64.8 & 100.00 & 0.00 & 0.00 & 0.00 \\
Spill stores & 81.9 & 99.93 & 0.07 & 0.07 & 0.00 \\
\multicolumn{2}{l}{\small\textit{L1H\%+L2H\%+DRAM\%=100\%}} \\
\bottomrule
\end{tabular}
\end{table}

\subsection{Train Input Results}

\textit{Note: this run used an older simulation without \texttt{-{}-l2cache}; cache counters are not available.}


\begin{table}[H]
\centering
\caption{ROI spill statistics --- \texttt{526.blender\_r} (train input)}
\label{tab:blender_train_roi}
\begin{tabular}{lr}
\toprule
\textbf{Metric} & \textbf{Value} \\
\midrule
ROI instructions & 721,231,642,103 \\
ROI loads        & 191,998,916,318 \\
ROI stores       & 29,533,390,237 \\
\textbf{ROI spills}   & \textbf{25,389,499,715} \\
\textbf{SP/inst\%}   & \textbf{3.52\%} \\
\textbf{SP/load\%}   & \textbf{13.22\%} \\
\bottomrule
\end{tabular}
\end{table}

\subsection{Ref Input Results}
\textit{Ref simulation in progress.}



%%──────────────────────────────────────────────────────────────────────────────
\section{\texttt{511.povray\_r}}
\label{sec:povray}

\paragraph{Benchmark summary.}
\texttt{511.povray\_r} is the SPEC CPU2017 rate variant of POV-Ray.
Ray-tracer with recursive intersection testing and lighting evaluation -- highest SP/inst\% in the suite.
The ROI is bounded by \texttt{m5\_work\_begin()} / \texttt{m5\_work\_end()}
around \texttt{StartRender() + cooperate loop}.

\subsection{Input Datasets}

\begin{table}[H]
\centering
\caption{Input dataset sizes for \texttt{511.povray\_r}}
\label{tab:povray_inputs}
\small
\begin{tabular}{lp{8cm}l}
\toprule
\textbf{Tier} & \textbf{Input} & \textbf{Relative scale} \\
\midrule
\textbf{test} & SPEC-benchmark-test.ini --- $50\times50$ pixel output & tiny render \\
train & SPEC-benchmark-train.ini --- ${\sim}200\times150$ pixels & ${\sim}12\times$ more pixels \\
ref & SPEC-benchmark-ref.ini --- $1280\times768$ pixels & $392\times$ more pixels \\
\bottomrule
\end{tabular}
\end{table}

\subsection{Test Input Results}

\begin{table}[H]
\centering
\caption{ROI spill statistics --- \texttt{511.povray\_r} (test input)}
\label{tab:povray_test_roi}
\begin{tabular}{lr}
\toprule
\textbf{Metric} & \textbf{Value} \\
\midrule
ROI instructions & 2,240,825,558 \\
ROI loads        & 716,248,306 \\
ROI stores       & 287,976,470 \\
\textbf{ROI spills}   & \textbf{180,043,488} \\
\textbf{SP/inst\%}   & \textbf{8.03\%} \\
\textbf{SP/load\%}   & \textbf{25.14\%} \\
Simulated time   & 3.978\,s \\
CPI              & 3.55 \\
\bottomrule
\end{tabular}
\end{table}

\begin{table}[H]
\centering
\caption{Spill cache hit/miss breakdown --- \texttt{511.povray\_r} (test input)}
\label{tab:povray_test_cache}
\begin{tabular}{lrrrrr}
\toprule
\textbf{} & \textbf{Total (M)} & \textbf{L1H\%} & \textbf{L1M\%} & \textbf{L2H\%} & \textbf{DRAM\%} \\
\midrule
Spill loads  & 180.0 & 99.90 & 0.10 & 0.10 & 0.00 \\
Spill stores & 200.3 & 99.65 & 0.35 & 0.35 & 0.00 \\
\multicolumn{2}{l}{\small\textit{L1H\%+L2H\%+DRAM\%=100\%}} \\
\bottomrule
\end{tabular}
\end{table}

\subsection{Train Input Results}

\textit{Note: this run used an older simulation without \texttt{-{}-l2cache}; cache counters are not available.}


\begin{table}[H]
\centering
\caption{ROI spill statistics --- \texttt{511.povray\_r} (train input)}
\label{tab:povray_train_roi}
\begin{tabular}{lr}
\toprule
\textbf{Metric} & \textbf{Value} \\
\midrule
ROI instructions & 28,619,546,491 \\
ROI loads        & 9,207,049,191 \\
ROI stores       & 3,677,218,308 \\
\textbf{ROI spills}   & \textbf{2,238,921,193} \\
\textbf{SP/inst\%}   & \textbf{7.82\%} \\
\textbf{SP/load\%}   & \textbf{24.32\%} \\
\bottomrule
\end{tabular}
\end{table}

\subsection{Ref Input Results}
\textit{Ref simulation in progress.}



%%──────────────────────────────────────────────────────────────────────────────
\section{\texttt{520.omnetpp\_r}}
\label{sec:omnetpp}

\paragraph{Benchmark summary.}
\texttt{520.omnetpp\_r} is the SPEC CPU2017 rate variant of OMNeT++.
Object-oriented discrete-event network simulator with deep virtual-dispatch call chains and many short-lived stack frames.
The ROI is bounded by \texttt{m5\_work\_begin()} / \texttt{m5\_work\_end()}
around \texttt{setupUserInterface()}.

\subsection{Input Datasets}

\begin{table}[H]
\centering
\caption{Input dataset sizes for \texttt{520.omnetpp\_r}}
\label{tab:omnetpp_inputs}
\small
\begin{tabular}{lp{8cm}l}
\toprule
\textbf{Tier} & \textbf{Input} & \textbf{Relative scale} \\
\midrule
\textbf{test} & General config, simulated time limit 0.003\,s & 0.003\,s sim \\
train & General config, simulated time limit 0.15\,s & 50$\times$ more sim time \\
ref & General config, simulated time limit 2.25\,s & 750$\times$ more sim time \\
\bottomrule
\end{tabular}
\end{table}

\subsection{Test Input Results}

\begin{table}[H]
\centering
\caption{ROI spill statistics --- \texttt{520.omnetpp\_r} (test input)}
\label{tab:omnetpp_test_roi}
\begin{tabular}{lr}
\toprule
\textbf{Metric} & \textbf{Value} \\
\midrule
ROI instructions & 10,286,636 \\
ROI loads        & 1,642,711 \\
ROI stores       & 655,212 \\
\textbf{ROI spills}   & \textbf{437,793} \\
\textbf{SP/inst\%}   & \textbf{4.26\%} \\
\textbf{SP/load\%}   & \textbf{26.65\%} \\
Simulated time   & 0.022\,s \\
CPI              & 4.37 \\
\bottomrule
\end{tabular}
\end{table}

\begin{table}[H]
\centering
\caption{Spill cache hit/miss breakdown --- \texttt{520.omnetpp\_r} (test input)}
\label{tab:omnetpp_test_cache}
\begin{tabular}{lrrrrr}
\toprule
\textbf{} & \textbf{Total (M)} & \textbf{L1H\%} & \textbf{L1M\%} & \textbf{L2H\%} & \textbf{DRAM\%} \\
\midrule
Spill loads  & 0.4 & 99.94 & 0.06 & 0.06 & 0.00 \\
Spill stores & 0.5 & 99.84 & 0.16 & 0.07 & 0.09 \\
\multicolumn{2}{l}{\small\textit{L1H\%+L2H\%+DRAM\%=100\%}} \\
\bottomrule
\end{tabular}
\end{table}

\subsection{Train Input Results}

\textit{Note: this run used an older simulation without \texttt{-{}-l2cache}; cache counters are not available.}


\begin{table}[H]
\centering
\caption{ROI spill statistics --- \texttt{520.omnetpp\_r} (train input)}
\label{tab:omnetpp_train_roi}
\begin{tabular}{lr}
\toprule
\textbf{Metric} & \textbf{Value} \\
\midrule
ROI instructions & 134,779,629,501 \\
ROI loads        & 34,419,513,685 \\
ROI stores       & 17,531,247,389 \\
\textbf{ROI spills}   & \textbf{9,219,009,244} \\
\textbf{SP/inst\%}   & \textbf{6.84\%} \\
\textbf{SP/load\%}   & \textbf{26.78\%} \\
\bottomrule
\end{tabular}
\end{table}

\subsection{Ref Input Results}
\textit{Ref simulation in progress.}



%%──────────────────────────────────────────────────────────────────────────────
\section{\texttt{523.xalancbmk\_r}}
\label{sec:xalancbmk}

\paragraph{Benchmark summary.}
\texttt{523.xalancbmk\_r} is the SPEC CPU2017 rate variant of Xalan-C.
XSLT document transformation; recursive XPath evaluation and heavy OO design drive persistent spilling.
The ROI is bounded by \texttt{m5\_work\_begin()} / \texttt{m5\_work\_end()}
around \texttt{xsltMain()}.

\subsection{Input Datasets}

\begin{table}[H]
\centering
\caption{Input dataset sizes for \texttt{523.xalancbmk\_r}}
\label{tab:xalancbmk_inputs}
\small
\begin{tabular}{lp{8cm}l}
\toprule
\textbf{Tier} & \textbf{Input} & \textbf{Relative scale} \\
\midrule
\textbf{test} & test.xml (28\,kB) + xalanc.xsl & small document \\
train & allbooks.xml (larger) + xalanc.xsl & larger document \\
ref & t5.xml (largest) + xalanc.xsl & largest document \\
\bottomrule
\end{tabular}
\end{table}

\subsection{Test Input Results}

\begin{table}[H]
\centering
\caption{ROI spill statistics --- \texttt{523.xalancbmk\_r} (test input)}
\label{tab:xalancbmk_test_roi}
\begin{tabular}{lr}
\toprule
\textbf{Metric} & \textbf{Value} \\
\midrule
ROI instructions & 395,919,950 \\
ROI loads        & 106,032,294 \\
ROI stores       & 38,990,477 \\
\textbf{ROI spills}   & \textbf{19,354,339} \\
\textbf{SP/inst\%}   & \textbf{4.89\%} \\
\textbf{SP/load\%}   & \textbf{18.25\%} \\
Simulated time   & 0.719\,s \\
CPI              & 3.63 \\
\bottomrule
\end{tabular}
\end{table}

\begin{table}[H]
\centering
\caption{Spill cache hit/miss breakdown --- \texttt{523.xalancbmk\_r} (test input)}
\label{tab:xalancbmk_test_cache}
\begin{tabular}{lrrrrr}
\toprule
\textbf{} & \textbf{Total (M)} & \textbf{L1H\%} & \textbf{L1M\%} & \textbf{L2H\%} & \textbf{DRAM\%} \\
\midrule
Spill loads  & 19.4 & 99.90 & 0.10 & 0.10 & 0.00 \\
Spill stores & 20.1 & 99.68 & 0.32 & 0.32 & 0.00 \\
\multicolumn{2}{l}{\small\textit{L1H\%+L2H\%+DRAM\%=100\%}} \\
\bottomrule
\end{tabular}
\end{table}

\subsection{Train Input Results}

\begin{table}[H]
\centering
\caption{ROI spill statistics --- \texttt{523.xalancbmk\_r} (train input)}
\label{tab:xalancbmk_train_roi}
\begin{tabular}{lr}
\toprule
\textbf{Metric} & \textbf{Value} \\
\midrule
ROI instructions & 7,848,338,475 \\
ROI loads        & 1,998,392,467 \\
ROI stores       & 956,498,570 \\
\textbf{ROI spills}   & \textbf{617,655,308} \\
\textbf{SP/inst\%}   & \textbf{7.87\%} \\
\textbf{SP/load\%}   & \textbf{30.91\%} \\
Simulated time   & 14.739\,s \\
CPI              & 3.76 \\
\bottomrule
\end{tabular}
\end{table}

\begin{table}[H]
\centering
\caption{Spill cache hit/miss breakdown --- \texttt{523.xalancbmk\_r} (train input)}
\label{tab:xalancbmk_train_cache}
\begin{tabular}{lrrrrr}
\toprule
\textbf{} & \textbf{Total (M)} & \textbf{L1H\%} & \textbf{L1M\%} & \textbf{L2H\%} & \textbf{DRAM\%} \\
\midrule
Spill loads  & 617.7 & 99.93 & 0.07 & 0.07 & 0.00 \\
Spill stores & 648.3 & 99.96 & 0.04 & 0.04 & 0.00 \\
\multicolumn{2}{l}{\small\textit{L1H\%+L2H\%+DRAM\%=100\%}} \\
\bottomrule
\end{tabular}
\end{table}

\subsection{Ref Input Results}
\textit{Ref simulation in progress.}



%%──────────────────────────────────────────────────────────────────────────────
\section{\texttt{525.x264\_r}}
\label{sec:x264}

\paragraph{Benchmark summary.}
\texttt{525.x264\_r} is the SPEC CPU2017 rate variant of x264/ldecod.
H.264 video decoder (ldecod path); inner decode loops maintain per-slice state that frequently spills.
The ROI is bounded by \texttt{m5\_work\_begin()} / \texttt{m5\_work\_end()}
around \texttt{do-while decode loop}.

\subsection{Input Datasets}

\begin{table}[H]
\centering
\caption{Input dataset sizes for \texttt{525.x264\_r}}
\label{tab:x264_inputs}
\small
\begin{tabular}{lp{8cm}l}
\toprule
\textbf{Tier} & \textbf{Input} & \textbf{Relative scale} \\
\midrule
\textbf{test} & BuckBunny\_train.264 --- 1.3\,MB H.264 stream (720p excerpt) & 1.3\,MB stream \\
train & BuckBunny.264 --- 2.4\,MB full H.264 file (720p) & 1.8$\times$ larger \\
ref & BuckBunny.264 --- 2.4\,MB full H.264 file (same as train) & same as train \\
\bottomrule
\end{tabular}
\end{table}

\subsection{Test Input Results}

\begin{table}[H]
\centering
\caption{ROI spill statistics --- \texttt{525.x264\_r} (test input)}
\label{tab:x264_test_roi}
\begin{tabular}{lr}
\toprule
\textbf{Metric} & \textbf{Value} \\
\midrule
ROI instructions & 18,192,367,319 \\
ROI loads        & 3,573,842,627 \\
ROI stores       & 1,995,653,850 \\
\textbf{ROI spills}   & \textbf{664,763,019} \\
\textbf{SP/inst\%}   & \textbf{3.65\%} \\
\textbf{SP/load\%}   & \textbf{18.60\%} \\
Simulated time   & 34.566\,s \\
CPI              & 3.80 \\
\bottomrule
\end{tabular}
\end{table}

\begin{table}[H]
\centering
\caption{Spill cache hit/miss breakdown --- \texttt{525.x264\_r} (test input)}
\label{tab:x264_test_cache}
\begin{tabular}{lrrrrr}
\toprule
\textbf{} & \textbf{Total (M)} & \textbf{L1H\%} & \textbf{L1M\%} & \textbf{L2H\%} & \textbf{DRAM\%} \\
\midrule
Spill loads  & 664.8 & 99.99 & 0.01 & 0.01 & 0.00 \\
Spill stores & 789.6 & 99.97 & 0.03 & 0.02 & 0.00 \\
\multicolumn{2}{l}{\small\textit{L1H\%+L2H\%+DRAM\%=100\%}} \\
\bottomrule
\end{tabular}
\end{table}

\subsection{Train Input Results}

\textit{Note: this run used an older simulation without \texttt{-{}-l2cache}; cache counters are not available.}


\begin{table}[H]
\centering
\caption{ROI spill statistics --- \texttt{525.x264\_r} (train input)}
\label{tab:x264_train_roi}
\begin{tabular}{lr}
\toprule
\textbf{Metric} & \textbf{Value} \\
\midrule
ROI instructions & 18,192,367,319 \\
ROI loads        & 3,573,842,627 \\
ROI stores       & 1,995,653,850 \\
\textbf{ROI spills}   & \textbf{638,746,003} \\
\textbf{SP/inst\%}   & \textbf{3.51\%} \\
\textbf{SP/load\%}   & \textbf{17.87\%} \\
\bottomrule
\end{tabular}
\end{table}

\subsection{Ref Input Results}

\begin{table}[H]
\centering
\caption{ROI spill statistics --- \texttt{525.x264\_r} (ref input)}
\label{tab:x264_ref_roi}
\begin{tabular}{lr}
\toprule
\textbf{Metric} & \textbf{Value} \\
\midrule
ROI instructions & 23,405,758,003 \\
ROI loads        & 4,837,558,149 \\
ROI stores       & 2,401,493,460 \\
\textbf{ROI spills}   & \textbf{804,721,799} \\
\textbf{SP/inst\%}   & \textbf{3.44\%} \\
\textbf{SP/load\%}   & \textbf{16.63\%} \\
Simulated time   & 43.901\,s \\
CPI              & 3.75 \\
\bottomrule
\end{tabular}
\end{table}

\begin{table}[H]
\centering
\caption{Spill cache hit/miss breakdown --- \texttt{525.x264\_r} (ref input)}
\label{tab:x264_ref_cache}
\begin{tabular}{lrrrrr}
\toprule
\textbf{} & \textbf{Total (M)} & \textbf{L1H\%} & \textbf{L1M\%} & \textbf{L2H\%} & \textbf{DRAM\%} \\
\midrule
Spill loads  & 804.7 & 99.99 & 0.01 & 0.01 & 0.00 \\
Spill stores & 952.3 & 99.97 & 0.03 & 0.03 & 0.00 \\
\multicolumn{2}{l}{\small\textit{L1H\%+L2H\%+DRAM\%=100\%}} \\
\bottomrule
\end{tabular}
\end{table}


%%──────────────────────────────────────────────────────────────────────────────
\section{\texttt{502.gcc\_r}}
\label{sec:gcc}

\paragraph{Benchmark summary.}
\texttt{502.gcc\_r} is the SPEC CPU2017 rate variant of GCC.
Full GCC compiler pass sequence; IR-manipulation passes keep many values simultaneously live, producing the highest SP/load\% in the suite.
The ROI is bounded by \texttt{m5\_work\_begin()} / \texttt{m5\_work\_end()}
around \texttt{toplev\_main()}.

\subsection{Input Datasets}

\begin{table}[H]
\centering
\caption{Input dataset sizes for \texttt{502.gcc\_r}}
\label{tab:gcc_inputs}
\small
\begin{tabular}{lp{8cm}l}
\toprule
\textbf{Tier} & \textbf{Input} & \textbf{Relative scale} \\
\midrule
\textbf{test} & t1.c --- a 69-byte trivial C source file, compiled at \texttt{-O3} & single tiny file \\
train & train01.c --- a 1.2\,MB C source file, compiled at \texttt{-O3} & ${\sim}17{,}000\times$ larger source file \\
ref & gcc-pp.c --- an 11\,MB C preprocessed source, compiled at \texttt{-O3} & ${\sim}159{,}000\times$ larger source file \\
\bottomrule
\end{tabular}
\end{table}

\subsection{Test Input Results}

\begin{table}[H]
\centering
\caption{ROI spill statistics --- \texttt{502.gcc\_r} (test input)}
\label{tab:gcc_test_roi}
\begin{tabular}{lr}
\toprule
\textbf{Metric} & \textbf{Value} \\
\midrule
ROI instructions & 16,027,045 \\
ROI loads        & 2,957,463 \\
ROI stores       & 1,955,576 \\
\textbf{ROI spills}   & \textbf{1,117,205} \\
\textbf{SP/inst\%}   & \textbf{6.97\%} \\
\textbf{SP/load\%}   & \textbf{37.78\%} \\
Simulated time   & 0.040\,s \\
CPI              & 4.93 \\
\bottomrule
\end{tabular}
\end{table}

\begin{table}[H]
\centering
\caption{Spill cache hit/miss breakdown --- \texttt{502.gcc\_r} (test input)}
\label{tab:gcc_test_cache}
\begin{tabular}{lrrrrr}
\toprule
\textbf{} & \textbf{Total (M)} & \textbf{L1H\%} & \textbf{L1M\%} & \textbf{L2H\%} & \textbf{DRAM\%} \\
\midrule
Spill loads  & 1.1 & 99.99 & 0.01 & 0.01 & 0.00 \\
Spill stores & 1.2 & 99.94 & 0.06 & 0.05 & 0.01 \\
\multicolumn{2}{l}{\small\textit{L1H\%+L2H\%+DRAM\%=100\%}} \\
\bottomrule
\end{tabular}
\end{table}

\subsection{Train Input Results}

\textit{Train simulation in progress.}


\subsection{Ref Input Results}

\begin{table}[H]
\centering
\caption{ROI spill statistics --- \texttt{502.gcc\_r} (ref input)}
\label{tab:gcc_ref_roi}
\begin{tabular}{lr}
\toprule
\textbf{Metric} & \textbf{Value} \\
\midrule
ROI instructions & 224,869,007,944 \\
ROI loads        & 54,829,285,747 \\
ROI stores       & 27,885,731,758 \\
\textbf{ROI spills}   & \textbf{18,182,291,057} \\
\textbf{SP/inst\%}   & \textbf{8.09\%} \\
\textbf{SP/load\%}   & \textbf{33.16\%} \\
Simulated time   & 455.531\,s \\
CPI              & 4.05 \\
\bottomrule
\end{tabular}
\end{table}

\begin{table}[H]
\centering
\caption{Spill cache hit/miss breakdown --- \texttt{502.gcc\_r} (ref input)}
\label{tab:gcc_ref_cache}
\begin{tabular}{lrrrrr}
\toprule
\textbf{} & \textbf{Total (M)} & \textbf{L1H\%} & \textbf{L1M\%} & \textbf{L2H\%} & \textbf{DRAM\%} \\
\midrule
Spill loads  & 18,182.3 & 99.99 & 0.01 & 0.01 & 0.00 \\
Spill stores & 19,612.9 & 99.79 & 0.21 & 0.21 & 0.00 \\
\multicolumn{2}{l}{\small\textit{L1H\%+L2H\%+DRAM\%=100\%}} \\
\bottomrule
\end{tabular}
\end{table}


%%──────────────────────────────────────────────────────────────────────────────
\section{\texttt{505.mcf\_r}}
\label{sec:mcf}

\paragraph{Benchmark summary.}
\texttt{505.mcf\_r} is the SPEC CPU2017 rate variant of MCF.
Min-cost flow network optimiser; sparse graph traversal with pointer-chasing limits register reuse.
The ROI is bounded by \texttt{m5\_work\_begin()} / \texttt{m5\_work\_end()}
around \texttt{global\_opt()}.

\subsection{Input Datasets}

\begin{table}[H]
\centering
\caption{Input dataset sizes for \texttt{505.mcf\_r}}
\label{tab:mcf_inputs}
\small
\begin{tabular}{lp{8cm}l}
\toprule
\textbf{Tier} & \textbf{Input} & \textbf{Relative scale} \\
\midrule
\textbf{test} & inp.in --- 1.2\,MB min-cost flow problem instance & 1.2\,MB instance \\
train & inp\_train.in --- 2.2\,MB larger problem instance & 1.8$\times$ larger \\
ref & inp.in --- same 1.2\,MB instance as test & identical to test \\
\bottomrule
\end{tabular}
\end{table}

\subsection{Test Input Results}

\begin{table}[H]
\centering
\caption{ROI spill statistics --- \texttt{505.mcf\_r} (test input)}
\label{tab:mcf_test_roi}
\begin{tabular}{lr}
\toprule
\textbf{Metric} & \textbf{Value} \\
\midrule
ROI instructions & 24,426,982,048 \\
ROI loads        & 8,113,038,194 \\
ROI stores       & 1,300,618,520 \\
\textbf{ROI spills}   & \textbf{509,449,025} \\
\textbf{SP/inst\%}   & \textbf{2.09\%} \\
\textbf{SP/load\%}   & \textbf{6.28\%} \\
Simulated time   & 93.459\,s \\
CPI              & 7.65 \\
\bottomrule
\end{tabular}
\end{table}

\begin{table}[H]
\centering
\caption{Spill cache hit/miss breakdown --- \texttt{505.mcf\_r} (test input)}
\label{tab:mcf_test_cache}
\begin{tabular}{lrrrrr}
\toprule
\textbf{} & \textbf{Total (M)} & \textbf{L1H\%} & \textbf{L1M\%} & \textbf{L2H\%} & \textbf{DRAM\%} \\
\midrule
Spill loads  & 509.4 & 99.88 & 0.12 & 0.12 & 0.00 \\
Spill stores & 705.2 & 99.54 & 0.46 & 0.46 & 0.01 \\
\multicolumn{2}{l}{\small\textit{L1H\%+L2H\%+DRAM\%=100\%}} \\
\bottomrule
\end{tabular}
\end{table}

\subsection{Train Input Results}

\textit{Note: this run used an older simulation without \texttt{-{}-l2cache}; cache counters are not available.}


\begin{table}[H]
\centering
\caption{ROI spill statistics --- \texttt{505.mcf\_r} (train input)}
\label{tab:mcf_train_roi}
\begin{tabular}{lr}
\toprule
\textbf{Metric} & \textbf{Value} \\
\midrule
ROI instructions & 125,387,999,080 \\
ROI loads        & 40,288,133,588 \\
ROI stores       & 8,306,287,899 \\
\textbf{ROI spills}   & \textbf{2,181,746,916} \\
\textbf{SP/inst\%}   & \textbf{1.74\%} \\
\textbf{SP/load\%}   & \textbf{5.42\%} \\
\bottomrule
\end{tabular}
\end{table}

\subsection{Ref Input Results}

\begin{table}[H]
\centering
\caption{ROI spill statistics --- \texttt{505.mcf\_r} (ref input)}
\label{tab:mcf_ref_roi}
\begin{tabular}{lr}
\toprule
\textbf{Metric} & \textbf{Value} \\
\midrule
ROI instructions & 24,426,982,048 \\
ROI loads        & 8,113,038,194 \\
ROI stores       & 1,300,618,520 \\
\textbf{ROI spills}   & \textbf{509,449,025} \\
\textbf{SP/inst\%}   & \textbf{2.09\%} \\
\textbf{SP/load\%}   & \textbf{6.28\%} \\
Simulated time   & 93.459\,s \\
CPI              & 7.65 \\
\bottomrule
\end{tabular}
\end{table}

\begin{table}[H]
\centering
\caption{Spill cache hit/miss breakdown --- \texttt{505.mcf\_r} (ref input)}
\label{tab:mcf_ref_cache}
\begin{tabular}{lrrrrr}
\toprule
\textbf{} & \textbf{Total (M)} & \textbf{L1H\%} & \textbf{L1M\%} & \textbf{L2H\%} & \textbf{DRAM\%} \\
\midrule
Spill loads  & 509.4 & 99.88 & 0.12 & 0.12 & 0.00 \\
Spill stores & 705.2 & 99.54 & 0.46 & 0.46 & 0.01 \\
\multicolumn{2}{l}{\small\textit{L1H\%+L2H\%+DRAM\%=100\%}} \\
\bottomrule
\end{tabular}
\end{table}


%%──────────────────────────────────────────────────────────────────────────────
\section{\texttt{500.perlbench\_r}}
\label{sec:perlbench}

\paragraph{Benchmark summary.}
\texttt{500.perlbench\_r} is the SPEC CPU2017 rate variant of Perlbench.
Perl interpreter; opcode dispatch loop and runtime maintain large per-op state, yielding high SP/inst\%.
The ROI is bounded by \texttt{m5\_work\_begin()} / \texttt{m5\_work\_end()}
around \texttt{perl\_run()}.

\subsection{Input Datasets}

\begin{table}[H]
\centering
\caption{Input dataset sizes for \texttt{500.perlbench\_r}}
\label{tab:perlbench_inputs}
\small
\begin{tabular}{lp{8cm}l}
\toprule
\textbf{Tier} & \textbf{Input} & \textbf{Relative scale} \\
\midrule
\textbf{test} & test.pl (29\,B) --- minimal Perl smoke-test script & trivial script \\
train & scrabbl.pl + suns.pl --- two real Perl programs & realistic workloads \\
ref & checkspam.pl / diffmail.pl / splitmail.pl (3 runs) & production scripts \\
\bottomrule
\end{tabular}
\end{table}

\subsection{Test Input Results}

\begin{table}[H]
\centering
\caption{ROI spill statistics --- \texttt{500.perlbench\_r} (test input)}
\label{tab:perlbench_test_roi}
\begin{tabular}{lr}
\toprule
\textbf{Metric} & \textbf{Value} \\
\midrule
ROI instructions & 14,690,864 \\
ROI loads        & 3,491,473 \\
ROI stores       & 2,035,472 \\
\textbf{ROI spills}   & \textbf{1,080,784} \\
\textbf{SP/inst\%}   & \textbf{7.36\%} \\
\textbf{SP/load\%}   & \textbf{30.95\%} \\
Simulated time   & 0.032\,s \\
CPI              & 4.34 \\
\bottomrule
\end{tabular}
\end{table}

\begin{table}[H]
\centering
\caption{Spill cache hit/miss breakdown --- \texttt{500.perlbench\_r} (test input)}
\label{tab:perlbench_test_cache}
\begin{tabular}{lrrrrr}
\toprule
\textbf{} & \textbf{Total (M)} & \textbf{L1H\%} & \textbf{L1M\%} & \textbf{L2H\%} & \textbf{DRAM\%} \\
\midrule
Spill loads  & 1.1 & 100.00 & 0.00 & 0.00 & 0.00 \\
Spill stores & 1.2 & 99.86 & 0.14 & 0.14 & 0.01 \\
\multicolumn{2}{l}{\small\textit{L1H\%+L2H\%+DRAM\%=100\%}} \\
\bottomrule
\end{tabular}
\end{table}

\subsection{Train Input Results}

\subsubsection*{Sub-run: \texttt{scrabbl.pl}}

\begin{table}[H]
\centering
\caption{ROI spill statistics --- \texttt{500.perlbench\_r} (train (scrabbl.pl) input)}
\label{tab:perl_scrabbl_train (scrabbl.pl)_roi}
\begin{tabular}{lr}
\toprule
\textbf{Metric} & \textbf{Value} \\
\midrule
ROI instructions & 1,176,588,509 \\
ROI loads        & 302,149,216 \\
ROI stores       & 182,574,177 \\
\textbf{ROI spills}   & \textbf{115,765,326} \\
\textbf{SP/inst\%}   & \textbf{9.84\%} \\
\textbf{SP/load\%}   & \textbf{38.31\%} \\
Simulated time   & 2.324\,s \\
CPI              & 3.95 \\
\bottomrule
\end{tabular}
\end{table}

\begin{table}[H]
\centering
\caption{Spill cache hit/miss breakdown --- \texttt{500.perlbench\_r} (train (scrabbl.pl) input)}
\label{tab:perl_scrabbl_train (scrabbl.pl)_cache}
\begin{tabular}{lrrrrr}
\toprule
\textbf{} & \textbf{Total (M)} & \textbf{L1H\%} & \textbf{L1M\%} & \textbf{L2H\%} & \textbf{DRAM\%} \\
\midrule
Spill loads  & 115.8 & 100.00 & 0.00 & 0.00 & 0.00 \\
Spill stores & 122.2 & 99.98 & 0.02 & 0.02 & 0.00 \\
\multicolumn{2}{l}{\small\textit{L1H\%+L2H\%+DRAM\%=100\%}} \\
\bottomrule
\end{tabular}
\end{table}

\subsubsection*{Sub-run: \texttt{suns.pl}}

\begin{table}[H]
\centering
\caption{ROI spill statistics --- \texttt{500.perlbench\_r} (train (suns.pl) input)}
\label{tab:perl_suns_train (suns.pl)_roi}
\begin{tabular}{lr}
\toprule
\textbf{Metric} & \textbf{Value} \\
\midrule
ROI instructions & 2,739,251,962 \\
ROI loads        & 681,242,110 \\
ROI stores       & 419,375,212 \\
\textbf{ROI spills}   & \textbf{285,586,649} \\
\textbf{SP/inst\%}   & \textbf{10.43\%} \\
\textbf{SP/load\%}   & \textbf{41.92\%} \\
Simulated time   & 4.646\,s \\
CPI              & 3.39 \\
\bottomrule
\end{tabular}
\end{table}

\begin{table}[H]
\centering
\caption{Spill cache hit/miss breakdown --- \texttt{500.perlbench\_r} (train (suns.pl) input)}
\label{tab:perl_suns_train (suns.pl)_cache}
\begin{tabular}{lrrrrr}
\toprule
\textbf{} & \textbf{Total (M)} & \textbf{L1H\%} & \textbf{L1M\%} & \textbf{L2H\%} & \textbf{DRAM\%} \\
\midrule
Spill loads  & 285.6 & 100.00 & 0.00 & 0.00 & 0.00 \\
Spill stores & 288.8 & 99.98 & 0.02 & 0.02 & 0.00 \\
\multicolumn{2}{l}{\small\textit{L1H\%+L2H\%+DRAM\%=100\%}} \\
\bottomrule
\end{tabular}
\end{table}

\subsection{Ref Input Results}
\textit{Ref simulation in progress.}



\clearpage
\section{Cross-Benchmark Summary}
\label{sec:summary}

The following subsections present summary tables for each simulation tier.
Benchmarks that do not yet have results for a given tier are omitted.
All tables are sorted by SP/inst\% (descending) to highlight the
benchmarks with the highest register pressure.


\clearpage
\subsection{Test Input --- Complete Results}
\label{subsec:summary_test}

Tables~\ref{tab:sum_test_spill} and~\ref{tab:sum_test_cache} list all
15 benchmarks sorted by \textbf{SP/inst\%} (descending).

\begin{landscape}
\begin{table}[p]
\centering
\caption{Test input: ROI spill statistics (sorted by SP/inst\%)}
\label{tab:sum_test_spill}
\small
\begin{tabular}{clrrrrrrr}
\toprule
\textbf{Rank} & \textbf{Benchmark} & \textbf{Insts (B)} & \textbf{Loads (B)} &
\textbf{Stores (B)} & \textbf{Spills (M)} & \textbf{SP/inst\%} &
\textbf{SP/load\%} & \textbf{CPI} \\
\midrule
1 & \texttt{511.povray\_r} & 2.24 & 0.72 & 0.29 & 180.0 & 8.03 & 25.14 & 3.55 \\
2 & \texttt{500.perlbench\_r} & 0.01 & 0.00 & 0.00 & 1.1 & 7.36 & 30.95 & 4.34 \\
3 & \texttt{502.gcc\_r} & 0.02 & 0.00 & 0.00 & 1.1 & 6.97 & 37.78 & 4.93 \\
4 & \texttt{526.blender\_r} & 0.99 & 0.23 & 0.12 & 64.8 & 6.52 & 27.75 & 4.18 \\
5 & \texttt{523.xalancbmk\_r} & 0.40 & 0.11 & 0.04 & 19.4 & 4.89 & 18.25 & 3.63 \\
6 & \texttt{541.leela\_r} & 26.06 & 5.11 & 1.76 & 1,179.4 & 4.52 & 23.08 & 3.33 \\
7 & \texttt{531.deepsjeng\_r} & 46.20 & 7.99 & 3.40 & 1,975.8 & 4.28 & 24.73 & 3.21 \\
8 & \texttt{520.omnetpp\_r} & 0.01 & 0.00 & 0.00 & 0.4 & 4.26 & 26.65 & 4.37 \\
9 & \texttt{519.lbm\_r} & 6.51 & 1.28 & 0.78 & 274.1 & 4.21 & 21.47 & 11.01 \\
10 & \texttt{525.x264\_r} & 18.19 & 3.57 & 2.00 & 664.8 & 3.65 & 18.60 & 3.80 \\
11 & \texttt{544.nab\_r} & 6.69 & 1.52 & 0.32 & 153.7 & 2.30 & 10.09 & 3.24 \\
12 & \texttt{505.mcf\_r} & 24.43 & 8.11 & 1.30 & 509.4 & 2.09 & 6.28 & 7.65 \\
13 & \texttt{538.imagick\_r} & 0.12 & 0.02 & 0.01 & 2.4 & 2.06 & 11.77 & 3.49 \\
14 & \texttt{557.xz\_r} & 2.01 & 0.37 & 0.18 & 39.2 & 1.95 & 10.65 & 4.77 \\
15 & \texttt{508.namd\_r} & 31.91 & 8.38 & 1.95 & 335.7 & 1.05 & 4.01 & 3.78 \\
\bottomrule
\end{tabular}
\end{table}

\begin{table}[p]
\centering
\caption{Test input: spill cache hit/miss breakdown (sorted by SP/inst\%, all \% of total)}
\label{tab:sum_test_cache}
\small
\begin{tabular}{clrrrrrrrrr}
\toprule
\multirow{2}{*}{\textbf{Rank}} & \multirow{2}{*}{\textbf{Benchmark}} &
\multicolumn{4}{c}{\textbf{Spill Loads}} & \phantom{x} &
\multicolumn{4}{c}{\textbf{Spill Stores}} \\
\cmidrule{3-6}\cmidrule{8-11}
 & & \textbf{L1H\%} & \textbf{L1M\%} & \textbf{L2H\%} & \textbf{DRAM\%} &
   & \textbf{L1H\%} & \textbf{L1M\%} & \textbf{L2H\%} & \textbf{DRAM\%} \\
\midrule
1 & \texttt{511.povray\_r} & 99.90 & 0.10 & 0.10 & 0.00 & & 99.65 & 0.35 & 0.35 & 0.00 \\
2 & \texttt{500.perlbench\_r} & 100.00 & 0.00 & 0.00 & 0.00 & & 99.86 & 0.14 & 0.14 & 0.01 \\
3 & \texttt{502.gcc\_r} & 99.99 & 0.01 & 0.01 & 0.00 & & 99.94 & 0.06 & 0.05 & 0.01 \\
4 & \texttt{526.blender\_r} & 100.00 & 0.00 & 0.00 & 0.00 & & 99.93 & 0.07 & 0.07 & 0.00 \\
5 & \texttt{523.xalancbmk\_r} & 99.90 & 0.10 & 0.10 & 0.00 & & 99.68 & 0.32 & 0.32 & 0.00 \\
6 & \texttt{541.leela\_r} & 100.00 & 0.00 & 0.00 & 0.00 & & 99.81 & 0.19 & 0.19 & 0.00 \\
7 & \texttt{531.deepsjeng\_r} & 99.97 & 0.03 & 0.03 & 0.00 & & 99.68 & 0.32 & 0.32 & 0.00 \\
8 & \texttt{520.omnetpp\_r} & 99.94 & 0.06 & 0.06 & 0.00 & & 99.84 & 0.16 & 0.07 & 0.09 \\
9 & \texttt{519.lbm\_r} & 100.00 & 0.00 & 0.00 & 0.00 & & 99.92 & 0.08 & 0.08 & 0.00 \\
10 & \texttt{525.x264\_r} & 99.99 & 0.01 & 0.01 & 0.00 & & 99.97 & 0.03 & 0.02 & 0.00 \\
11 & \texttt{544.nab\_r} & 100.00 & 0.00 & 0.00 & 0.00 & & 99.99 & 0.01 & 0.01 & 0.00 \\
12 & \texttt{505.mcf\_r} & 99.88 & 0.12 & 0.12 & 0.00 & & 99.54 & 0.46 & 0.46 & 0.01 \\
13 & \texttt{538.imagick\_r} & 100.00 & 0.00 & 0.00 & 0.00 & & 99.74 & 0.26 & 0.22 & 0.04 \\
14 & \texttt{557.xz\_r} & 100.00 & 0.00 & 0.00 & 0.00 & & 99.90 & 0.10 & 0.10 & 0.00 \\
15 & \texttt{508.namd\_r} & 99.93 & 0.07 & 0.07 & 0.00 & & 99.82 & 0.18 & 0.18 & 0.00 \\
\bottomrule
\end{tabular}
\end{table}
\end{landscape}


\subsection{Train Input --- Available Results}
\label{subsec:summary_train}

Table~\ref{tab:sum_train_spill} lists all benchmarks for which train-input
simulations are complete.  Entries marked with a dagger~($\dagger$) come from
runs \emph{without} \texttt{--l2cache} (no cache counters available).
Table~\ref{tab:sum_train_cache} shows cache breakdown for the subset run with
\texttt{--l2cache}.

\begin{table}[H]
\centering
\caption{Train input: ROI spill statistics (sorted by SP/inst\%)}
\label{tab:sum_train_spill}
\small
\begin{tabular}{clrrrr}
\toprule
\textbf{Rank} & \textbf{Benchmark} & \textbf{Insts (B)} & \textbf{Spills (M)}
  & \textbf{SP/inst\%} & \textbf{SP/load\%} \\
\midrule
1 & \texttt{500.perlbench\_r} \textit{(suns.pl)} & 2.74 & 285.6 & 10.43 & 41.92 \\
2 & \texttt{500.perlbench\_r} \textit{(scrabbl.pl)} & 1.18 & 115.8 & 9.84 & 38.31 \\
3 & \texttt{523.xalancbmk\_r} & 7.85 & 617.7 & 7.87 & 30.91 \\
4 & \texttt{511.povray\_r}$^\dagger$ & 28.62 & 2,238.9 & 7.82 & 24.32 \\
5 & \texttt{520.omnetpp\_r}$^\dagger$ & 134.78 & 9,219.0 & 6.84 & 26.78 \\
6 & \texttt{531.deepsjeng\_r}$^\dagger$ & 399.99 & 18,131.5 & 4.53 & 25.05 \\
7 & \texttt{541.leela\_r}$^\dagger$ & 436.32 & 19,460.7 & 4.46 & 22.81 \\
8 & \texttt{519.lbm\_r}$^\dagger$ & 100.16 & 4,229.6 & 4.22 & 21.61 \\
9 & \texttt{526.blender\_r}$^\dagger$ & 721.23 & 25,389.5 & 3.52 & 13.22 \\
10 & \texttt{525.x264\_r}$^\dagger$ & 18.19 & 638.7 & 3.51 & 17.87 \\
11 & \texttt{557.xz\_r}$^\dagger$ & 45.07 & 1,422.4 & 3.16 & 16.20 \\
12 & \texttt{505.mcf\_r}$^\dagger$ & 125.39 & 2,181.7 & 1.74 & 5.42 \\
13 & \texttt{508.namd\_r}$^\dagger$ & 218.43 & 2,285.7 & 1.05 & 3.96 \\
14 & \texttt{538.imagick\_r}$^\dagger$ & 284.11 & 292.4 & 0.10 & 0.61 \\
\bottomrule
\end{tabular}
\end{table}

\begin{table}[H]
\centering
\caption{Train input: spill cache hit/miss breakdown (\texttt{--l2cache} runs, all \% of total)}
\label{tab:sum_train_cache}
\small
\begin{tabular}{lrrrrrrrrrr}
\toprule
\multirow{2}{*}{\textbf{Benchmark}} &
\multicolumn{4}{c}{\textbf{Spill Loads}} & \phantom{x} &
\multicolumn{4}{c}{\textbf{Spill Stores}} \\
\cmidrule{2-5}\cmidrule{7-10}
 & \textbf{L1H\%} & \textbf{L1M\%} & \textbf{L2H\%} & \textbf{DRAM\%} &
   & \textbf{L1H\%} & \textbf{L1M\%} & \textbf{L2H\%} & \textbf{DRAM\%} \\
\midrule
\texttt{500.perlbench\_r} \textit{(suns.pl)} & 100.00 & 0.00 & 0.00 & 0.00 & & 99.98 & 0.02 & 0.02 & 0.00 \\
\texttt{500.perlbench\_r} \textit{(scrabbl.pl)} & 100.00 & 0.00 & 0.00 & 0.00 & & 99.98 & 0.02 & 0.02 & 0.00 \\
\texttt{523.xalancbmk\_r} & 99.93 & 0.07 & 0.07 & 0.00 & & 99.96 & 0.04 & 0.04 & 0.00 \\
\bottomrule
\end{tabular}
\end{table}


\subsection{Ref Input --- Available Results}
\label{subsec:summary_ref}

Table~\ref{tab:sum_ref_spill} lists benchmarks for which ref-input
simulations are complete.  Four benchmarks (\texttt{gcc}, \texttt{x264},
\texttt{mcf}, \texttt{imagick}) were run with \texttt{--l2cache} and
provide full cache statistics; \texttt{xz} provides spill counts only.

\begin{table}[H]
\centering
\caption{Ref input: ROI spill statistics (sorted by SP/inst\%)}
\label{tab:sum_ref_spill}
\small
\begin{tabular}{clrrrr}
\toprule
\textbf{Rank} & \textbf{Benchmark} & \textbf{Insts (B)} & \textbf{Spills (M)}
  & \textbf{SP/inst\%} & \textbf{SP/load\%} \\
\midrule
1 & \texttt{502.gcc\_r} & 224.87 & 18,182.3 & 8.09 & 33.16 \\
2 & \texttt{525.x264\_r} & 23.41 & 804.7 & 3.44 & 16.63 \\
3 & \texttt{505.mcf\_r} & 24.43 & 509.4 & 2.09 & 6.28 \\
4 & \texttt{538.imagick\_r} & 0.12 & 2.4 & 2.06 & 11.77 \\
5 & \texttt{557.xz\_r} & 401.33 & 6,636.2 & 1.65 & 10.15 \\
\bottomrule
\end{tabular}
\end{table}

\begin{table}[H]
\centering
\caption{Ref input: spill cache hit/miss breakdown (benchmarks with \texttt{--l2cache}, all \% of total)}
\label{tab:sum_ref_cache}
\small
\begin{tabular}{lrrrrrrrrrr}
\toprule
\multirow{2}{*}{\textbf{Benchmark}} &
\multicolumn{4}{c}{\textbf{Spill Loads}} & \phantom{x} &
\multicolumn{4}{c}{\textbf{Spill Stores}} \\
\cmidrule{2-5}\cmidrule{7-10}
 & \textbf{L1H\%} & \textbf{L1M\%} & \textbf{L2H\%} & \textbf{DRAM\%} &
   & \textbf{L1H\%} & \textbf{L1M\%} & \textbf{L2H\%} & \textbf{DRAM\%} \\
\midrule
\texttt{502.gcc\_r} & 99.99 & 0.01 & 0.01 & 0.00 & & 99.79 & 0.21 & 0.21 & 0.00 \\
\texttt{525.x264\_r} & 99.99 & 0.01 & 0.01 & 0.00 & & 99.97 & 0.03 & 0.03 & 0.00 \\
\texttt{505.mcf\_r} & 99.88 & 0.12 & 0.12 & 0.00 & & 99.54 & 0.46 & 0.46 & 0.01 \\
\texttt{538.imagick\_r} & 100.00 & 0.00 & 0.00 & 0.00 & & 99.74 & 0.26 & 0.22 & 0.04 \\
\bottomrule
\end{tabular}
\end{table}


\subsection{Discussion}

\paragraph{Spill rate range.}
SP/inst\% spans an $8\times$ range (1.05\%--8.03\%) across the test
tier.  The highest spilling benchmarks share a common trait: they
execute deep recursive call graphs that keep many values simultaneously
live.  \texttt{povray} (ray-tracer), \texttt{perlbench} (Perl
interpreter), \texttt{gcc} (compiler IR passes), and \texttt{blender}
(path-tracer) all exceed 6.5\%.

\paragraph{Test vs.\ train scaling.}
For benchmarks where both tiers are available, SP/inst\% is
remarkably stable.  \texttt{lbm} goes from 4.21\% (test, 20 steps)
to 4.22\% (train, 300 steps)---essentially constant---because each
time-step executes the identical D3Q19 collision kernel.
\texttt{deepsjeng} similarly holds near 4.5\% across input sizes.
\texttt{imagick} is a notable exception: 2.06\% (test, 318\,B input)
vs.\ 0.10\% (train, 284\,B input with 8 filters),
reflecting a qualitative difference in which code paths dominate for
the two filter chains.

\paragraph{SP/load\% as the key overhead metric.}
\texttt{gcc} is the clearest example: its 37.78\% SP/load\% means
that eliminating all register spilling would reduce its load traffic
by more than a third---a potential CPI improvement of similar
magnitude on a bandwidth-constrained pipeline.

\paragraph{Cache effectiveness.}
The L1-D cache absorbs $\geq$99.88\% of spill loads (L1H\%$\geq$99.88\%)
across all benchmarks and all tiers with cache data.  This holds even for
\texttt{gcc} ref (18 billion spills): the 8.09\% SP/inst\% leads to
just 1.6M L1 misses out of 18.2B total spill loads.
The conclusion is consistent: \emph{spilling costs instruction
bandwidth, not memory bandwidth}.

\paragraph{DRAM traffic from spilling.}
DRAML\% is 0.00\% for every benchmark across all tiers with cache data,
meaning no spill-load misses escape to DRAM.  L2H\%$\approx$L1M\%
for each benchmark, confirming that L2 rescues every L1 spill miss.
DRAM traffic from spilling is therefore zero or negligible ($<$0.1\%).

\section{Follow-Up Experiment: Small-Cache Robustness Check (\lowercase{v2} Spill Detector)}
\label{sec:smallcache}

The results above (Section~\ref{subsec:summary_test}, Discussion) were
obtained with generous cache sizes (L1-D 32\,kB 8-way, L2 256\,kB) and
showed DRAM traffic from spilling at 0.00\%.  A natural objection is
that this null result could simply be an artifact of over-provisioned
caches: with 32\,kB of L1-D alone, almost any per-call-frame spill
footprint would fit.  This section describes a follow-up experiment
designed specifically to stress-test that objection.

\subsection{Motivation}

If the near-zero DRAM traffic reported in Table~\ref{tab:sum_test_cache}
were a sizing artifact, then aggressively shrinking the cache hierarchy
should reveal spill working sets that no longer fit in L1/L2, pushing
measurable traffic to DRAM. If, instead, spill store/reload pairs have
inherently short reuse distances (bounded by the lifetime of a single
stack frame), then even a small cache should continue to absorb them
almost entirely --- confirming that the original conclusion is not an
artifact of cache sizing.

\subsection{Configuration Changes}

Two independent changes were introduced relative to the results in
Sections~\ref{sec:deepsjeng}--\ref{subsec:summary_train}:

\begin{table}[H]
\centering
\caption{Small-cache configuration vs.\ original}
\label{tab:smallcache_config}
\begin{tabular}{lrrr}
\toprule
\textbf{Component} & \textbf{Original} & \textbf{Small-cache} & \textbf{Reduction} \\
\midrule
L1-D & 32\,kB (8-way) & 4\,kB & $16\times$ \\
L1-I & 32\,kB & 32\,kB & unchanged \\
L2 (unified) & 256\,kB (8-way) & 32\,kB & $8\times$ \\
\bottomrule
\end{tabular}
\end{table}

\texttt{gem5} flags: \texttt{--l1d\_size=4kB --l1i\_size=32kB
--l2\_size=32kB} added to the base \texttt{--caches --l2cache}
configuration used throughout this report.

Separately, the \texttt{SpillDetector} module (Section~\ref{sec:method})
was revised (``v2''):
\begin{itemize}
  \item The original 50-million-tick staleness window on pending store
        entries was removed --- a spill pair is now matched regardless
        of the elapsed time between store and reload.
  \item A store/reload pair is only counted as a genuine spill if
        \emph{both} the store and the reload use SP/FP-relative
        addressing; pairs that satisfy this on only one side are stack
        data reuse through a computed pointer (arrays, structs) and are
        now excluded, tightening the definition.
  \item A spilled slot can be reloaded multiple times before the next
        store overwrites it; every such reload is now counted
        individually in \texttt{roi\_spills} (and tagged for cache
        statistics), not just the first.
\end{itemize}

\subsection{Results}

Six of the eight benchmarks selected for this follow-up
(run-time budget $\leq$20\,h under the original cache config) have
completed under the small-cache v2 configuration with the \textbf{train}
input tier: \texttt{xalancbmk}, \texttt{gcc}, \texttt{povray},
\texttt{x264}, \texttt{nab}, \texttt{xz}. \texttt{mcf} and
\texttt{omnetpp} were excluded from the core batch (estimated
30--35\,h each) and are pending as a separate tier-2 run; they are not
included in the numbers below.

\begin{table}[H]
\centering
\caption{Small-cache v2, train input: ROI spill statistics (sorted by SP/inst\%)}
\label{tab:sc_v2_spill}
\small
\begin{tabular}{clrrrr}
\toprule
\textbf{Rank} & \textbf{Benchmark} & \textbf{Insts (B)} & \textbf{Loads (B)} &
\textbf{SP/inst\%} & \textbf{SP/load\%} \\
\midrule
1 & \texttt{502.gcc\_r}       & 10.03 & 2.51 & 7.94 & 31.76 \\
2 & \texttt{523.xalancbmk\_r} &  7.85 & 2.00 & 7.71 & 30.27 \\
3 & \texttt{511.povray\_r}    & 28.62 & 9.21 & 7.58 & 23.55 \\
4 & \texttt{525.x264\_r}      & 23.41 & 4.84 & 3.14 & 15.19 \\
5 & \texttt{557.xz\_r}        & 45.07 & 8.78 & 2.82 & 14.48 \\
6 & \texttt{544.nab\_r}       & 38.23 & 8.73 & 2.03 &  8.89 \\
\bottomrule
\end{tabular}
\end{table}

\begin{table}[H]
\centering
\caption{Small-cache v2, train input: spill cache hit/miss breakdown (all \% of total)}
\label{tab:sc_v2_cache}
\small
\begin{tabular}{clrrrrrrrrr}
\toprule
\multirow{2}{*}{\textbf{Rank}} & \multirow{2}{*}{\textbf{Benchmark}} &
\multicolumn{4}{c}{\textbf{Spill Loads}} & \phantom{x} &
\multicolumn{4}{c}{\textbf{Spill Stores}} \\
\cmidrule{3-6}\cmidrule{8-11}
 & & \textbf{L1H\%} & \textbf{L1M\%} & \textbf{L2H\%} & \textbf{DRAM\%} &
   & \textbf{L1H\%} & \textbf{L1M\%} & \textbf{L2H\%} & \textbf{DRAM\%} \\
\midrule
1 & \texttt{502.gcc\_r}       & 98.89 & 1.11 & 1.08 & 0.04 & & 97.62 & 2.38 & 2.17 & 0.21 \\
2 & \texttt{523.xalancbmk\_r} & 98.06 & 1.94 & 1.79 & 0.15 & & 96.57 & 3.43 & 3.20 & 0.23 \\
3 & \texttt{511.povray\_r}    & 97.98 & 2.02 & 1.55 & 0.47 & & 97.62 & 2.38 & 1.37 & 1.01 \\
4 & \texttt{525.x264\_r}      & 97.69 & 2.31 & 2.30 & 0.01 & & 97.56 & 2.44 & 2.29 & 0.15 \\
5 & \texttt{557.xz\_r}        & 98.28 & 1.72 & 1.72 & 0.01 & & 95.20 & 4.80 & 4.78 & 0.02 \\
6 & \texttt{544.nab\_r}       & 98.95 & 1.05 & 1.05 & 0.00 & & 99.27 & 0.73 & 0.72 & 0.01 \\
\bottomrule
\end{tabular}
\end{table}

\subsection{Discussion}

\paragraph{The objection does not hold.} Even after cutting L1-D
$16\times$ and L2 $8\times$, L1H\% for spill loads remains
$\geq$97.69\% (vs.\ $\geq$99.88\% under the original config) and
DRAM\% stays $\leq$0.47\% for loads and $\leq$1.01\% for stores ---
\texttt{povray} in both cases, the benchmark with the largest ROI
in this set (28.6B instructions, 2.17B spills). Every other benchmark
stays at or below 0.23\% DRAM traffic on both loads and stores. The
order-of-magnitude increase in miss rate that a pure capacity argument
would predict does not appear.

\paragraph{Interpretation.} L2H\%$\approx$L1M\% continues to hold for
every benchmark (e.g.\ \texttt{xz}: L1M\%=4.80, L2H\%=4.78 for
stores), meaning L2 --- now only 32\,kB --- still rescues essentially
every L1 miss. This is consistent with spill store/reload pairs having
a short reuse distance bounded by the enclosing stack frame's
lifetime: the live spill footprint of a single function invocation is
small enough (tens to a few hundred bytes for these benchmarks) that
even a heavily undersized cache holds it. \textbf{The low DRAM traffic
observed in Section~\ref{sec:method} is therefore not a cache-sizing
artifact} --- it is a property of how register spilling is used by the
compiler's code generator, and it survives a $16\times$/$8\times$
reduction in L1/L2 capacity together with a stricter (v2) spill
definition.

\paragraph{Open work.} The two highest-instruction-count benchmarks in
the original suite, \texttt{mcf} (pointer-chasing, poor spatial
locality) and \texttt{omnetpp} (deep virtual-dispatch call chains),
are the most plausible candidates to break this pattern and are
queued as a tier-2 run (Section~\ref{sec:smallcache}, Results). Their
estimated 30--35\,h run-time exceeds the budget available at the time
of writing.

\section{Conclusion}

All 15 SPEC CPU2017 C/C++ benchmarks exhibit measurable register
spilling on RISC-V RV64GC.  Spill densities range 1--8\% of
instructions and 4--38\% of loads.  Across all simulation tiers
and all benchmarks for which cache data is available:
\begin{itemize}
  \item The L1-D cache serves $>$99.88\% of spill reloads (L1H\%$>$99.88\%).
  \item L2 rescues every L1 spill miss (DRAML\%$=$0.00\%).
  \item DRAM traffic from spilling is $<$0.1\% or zero (DRAMS\%$\leq$0.09\%).
\end{itemize}
The dominant cost of register spilling on this architecture is
therefore \emph{instruction-stream overhead} (added load/store
instructions and issue slots), not off-chip bandwidth.
\texttt{gcc} and \texttt{perlbench} are the most likely candidates
to benefit from compiler improvements or a hardware register-file
extension.

A follow-up robustness check (Section~\ref{sec:smallcache}) confirms
this holds even under a $16\times$/$8\times$ shrink of L1-D/L2 and a
stricter spill definition: DRAM traffic from spilling stays
$\leq$1.01\% for six of eight tested benchmarks, ruling out
over-provisioned caches as the explanation. \texttt{mcf} and
\texttt{omnetpp} remain to be measured under this configuration.

\end{document}
